\chapter{Исследовательская часть}

\section{Цель исследования}

При работе информационной системы поиска игровых модификаций значительная часть нагрузки приходится на выполнение поисковых запросов. По мере роста объема данных возрастает время поиска, что негативно влияет на производительность системы и пользовательский опыт.

Одним из основных механизмов повышения производительности является использование индексов. Однако различные типы индексов ориентированы на разные сценарии использования и могут демонстрировать различные результаты в зависимости от характера запросов.

Целью исследования является анализ влияния различных типов индексов PostgreSQL на скорость поиска игровых модификаций.

\section{Описание исследования}

Для проведения исследования была подготовлена база данных, содержащая информацию о игровых модификациях.

Были рассмотрены следующие варианты организации хранения данных:

\begin{itemize}
	\item отсутствие индекса;
	\item индекс B-Tree;
	\item индекс Hash;
	\item индекс GIN;
	\item индекс GiST;
	\item индекс BRIN.
\end{itemize}

Была составлена база данных, состоящая из 100\,000 различных записей модификаций.
Для каждого варианта индекса выполнялись следующие запросы:

\begin{itemize}
	\item поиск модификации по точному совпадению значения (названию);
	\item поиск модификаций в заданном диапазоне значений (дата публикации);
\end{itemize}

Для повышения достоверности результатов выполнялось несколько запусков каждого запроса с последующим вычислением среднего значения.

\section{Сравнение размеров индексов}

Размеры созданных индексов представлены на рисунке~\ref{fig:index_size}.

\begin{figure}[H]
	\centering
	\includegraphics[width=0.9\textwidth]{index_size}
	\caption{Размеры исследуемых индексов}
	\label{fig:index_size}
\end{figure}

Наименьший размер продемонстрировал индекс BRIN, занимающий всего 24 КБ. Размер индекса B-Tree составил 745 КБ. Индексы Hash, GIN и GiST потребовали значительно больше дискового пространства: 5.8 МБ, 6.1 МБ и 5.3 МБ соответственно.

Таким образом, по объему занимаемой памяти наименьшим является индекс BRIN, тогда как индексы GIN и Hash требуют более чем в 200 раз больше памяти.

\newpage
\section{Исследование скорости поиска по точному совпадению}

Результаты измерений представлены на рисунке~\ref{fig:equality_search}.

\begin{figure}[H]
	\centering
	\includegraphics[width=0.9\textwidth]{equality_search_time}
	\caption{Время поиска по точному совпадению}
	\label{fig:equality_search}
\end{figure}

Результаты исследования представлены в таблице~\ref{tab:equality}.

\begin{table}[H]
	\centering
	\caption{Время поиска по точному совпадению}
	\label{tab:equality}
	\begin{tabularx}{\textwidth}{|X|X|}
		\hline
		Тип индекса & Время, нс \\
		\hline
		Без индекса & 9\,596\,800 \\
		B-Tree & 3\,600\,858 \\
		Hash & 3\,098\,758 \\
		GIN & 9\,009\,500 \\
		GiST & 9\,267\,383 \\
		BRIN & 9\,275\,283 \\
		\hline
	\end{tabularx}
\end{table}

Наилучший результат продемонстрировал индекс Hash, обеспечивший выполнение запроса за 3.1 миллисекунды. По сравнению с отсутствием индекса ускорение составило примерно 68\%.

Индекс B-Tree показал сопоставимый результат --- 3.6 миллисекунды, что соответствует ускорению на 62.5\%.

Индексы GIN, GiST и BRIN практически не повлияли на производительность поиска по точному совпадению. Время выполнения запросов оказалось близким к варианту без индекса.

Полученные результаты объясняются тем, что Hash и B-Tree специально оптимизированы для операций сравнения по равенству.

\section{Исследование диапазонного поиска}

Результаты диапазонного поиска представлены на рисунке~\ref{fig:range_search}.

\begin{figure}[H]
	\centering
	\includegraphics[width=0.9\textwidth]{range_search_time}
	\caption{Время диапазонного поиска}
	\label{fig:range_search}
\end{figure}

Результаты измерений приведены в таблице~\ref{tab:range_search}.

\begin{table}[H]
	\centering
	\caption{Время диапазонного поиска}
	\label{tab:range_search}
	\begin{tabularx}{\textwidth}{|X|X|}
		\hline
		Тип индекса & Время, нс \\
		\hline
		Без индекса & 6\,159\,675 \\
		B-Tree & 6\,036\,067 \\
		Hash & 6\,354\,733 \\
		GIN & 6\,099\,567 \\
		GiST & 5\,931\,292 \\
		BRIN & 392\,342 \\
		\hline
	\end{tabularx}
\end{table}

Самым быстрым оказался индекс BRIN. Время выполнения запроса составило всего 392 тыс. нс, что в 15.7 раза быстрее последовательного сканирования таблицы.

Остальные типы индексов показали результаты, близкие друг к другу. Улучшение производительности относительно отсутствия индекса оказалось незначительным.

Высокая скорость BRIN объясняется тем, что данный индекс хранит сведения о диапазонах значений и позволяет быстро исключать большие фрагменты таблицы, не содержащие подходящих записей.

\section{Вывод}

Проведенное исследование показало, что скорость поиска при наличии индексов PostgreSQL существенно зависит от характера выполняемых запросов.

Для поиска игровых модификаций по точному совпадению наилучшие результаты продемонстрировали индексы Hash и B-Tree. Минимальное время выполнения запроса было получено при использовании Hash-индекса и составило 3.1 миллисекунды, что примерно в 3.1 раза быстрее по сравнению с поиском без индекса. Однако размер Hash-индекса составил 5.8 МБ, что почти в восемь раз превышает размер индекса B-Tree. Индекс B-Tree показал близкий результат --- 3.6 миллисекунды, обеспечив ускорение поиска более чем в 2.5 раза при значительно меньших затратах памяти. Это позволяет рассматривать B-Tree как лучшее решение для реализации поиска модификаций по наименованию среди других индексов.

При выполнении диапазонных запросов наилучший результат продемонстрировал индекс BRIN. Время выполнения запроса составило 392 тыс. нс, что в 15.7 раза быстрее по сравнению с последовательным сканированием таблицы. Дополнительным преимуществом данного индекса является минимальный объем занимаемой памяти --- 24 КБ. Полученные результаты подтверждают высокую скорость BRIN для работы с большими объемами данных, имеющими естественную упорядоченность значений.

Таким образом, результаты исследования позволяют сделать вывод, что для разработанной базы данных поиска игровых модификаций наиболее рациональным решением является использование индекса B-Tree в качестве основного механизма индексирования. Данный тип индекса обеспечивает высокую скорость поиска при умеренных затратах памяти и является наиболее универсальным для выполнения типовых запросов. Для операций диапазонного поиска целесообразно дополнительно использовать индекс BRIN, который обеспечивает максимальную скорость поиска при минимальном размере индекса.
